\chapter{Swap-Spread and Asset-Swap Trades}\label{ch:s2:swap-spread-and-asset-swap-trades}

From October 2008 to October 2016 the thirty-year US dollar swap rate stood below the Treasury yield of the same maturity on 94\% of days. A textbook
says this cannot last: buy the Treasury, pay fixed in the swap, and collect the difference. It lasted because the trade that would close it uses
balance sheet, and banks were no longer willing to lend theirs cheaply. Jermann showed that when holding bonds carries frictions, negative swap
spreads should not surprise anyone. On this chapter's synthetic market the long-swap-spread trade earns 32.4 basis points a year on its notional once
a \omterm{def:s2:swap-spread-and-asset-swap-trades:bscost}{balance-sheet cost} has pushed the spread negative. Charged for the balance sheet it uses, it loses 17.6. The build is \texttt{firm.swapspread}.

\section{Swap spreads and asset swaps}

A swap spread (Book~2, chapter~9) is the fixed rate of an interest rate swap minus the government bond yield of the same maturity. An asset-swap
spread (Book~2, chapter~21) is the spread over the floating rate that a bond pays once swapped: buying the bond and paying fixed on a swap with the
bond's coupons turns it into a floating-rate asset. Before 2008 swap spreads were positive, a premium for the bank credit in the floating leg and for
Treasuries' liquidity. From July 2000 to September 2008 the ten-year spread averaged 58.6 basis points and the thirty-year 48.3.

\begin{definition}[Swap-spread trade]\label{def:s2:swap-spread-and-asset-swap-trades:trade}
A \emph{swap-spread trade}\index{swap-spread trade} combines a government bond with an interest rate swap of the same maturity so that the position
gains when the swap spread moves: long the spread is long the bond (financed in repo) and paying fixed in the swap, which earns the bond yield over the
swap rate plus the floating rate over repo, and gains when the spread widens.
\end{definition}

\section{Negative swap spreads}

From October 2008 the long end turned negative (\cref{fig:s2:swap-spread-and-asset-swap-trades:real}). The thirty-year spread averaged $-21.4$ basis
points to October 2016, was negative on 93.9\% of those days, and reached $-62$ on 13 September 2016. The ten-year spread averaged 8.3 and was negative
on 17.6\% of days, reaching $-23$ on 24 February 2016. The two-year spread never went below zero: its minimum was 2 basis points, on 20 November 2015.
Jermann's model prices swaps when holding bonds is costly. It matches the negative spreads without large demand imbalances in the swap market.

\begin{omfigure}
\begin{tikzpicture}
\begin{axis}[width=11cm,height=5.4cm,xlabel={\small year},ylabel={\small swap $-$ Treasury (bp)},tick label style={font=\footnotesize},
  xmin=2000.5,xmax=2017,ymin=-70,ymax=140,grid=major,grid style={gray!25},table/col sep=comma,/pgf/number format/1000 sep={},
  legend style={font=\scriptsize,at={(0.97,0.97)},anchor=north east}]
\addplot[omDef,thick] table[x=year,y=s10]{figdata/strategies-2/12-swap-spread-and-asset-swap-trades/real.csv};
\addplot[omThm,thick] table[x=year,y=s30]{figdata/strategies-2/12-swap-spread-and-asset-swap-trades/real.csv};
\addplot[black,thin,domain=2000.5:2017] {0};
\legend{10 years,30 years}
\end{axis}
\end{tikzpicture}

\omcaption{US dollar swap spreads (the ICE swap rate minus the constant-maturity Treasury yield), monthly means, July 2000 to October 2016, when the
swap series on FRED end. Derived from FRED; the swap rates are not redistributed. Data: \texttt{s2\_fetch\_swaps}.}\label{fig:s2:swap-spread-and-asset-swap-trades:real}
\end{omfigure}

\section{Balance sheet as a cost}

\begin{definition}[Balance-sheet cost]\label{def:s2:swap-spread-and-asset-swap-trades:bscost}
The \emph{balance-sheet cost}\index{balance-sheet cost} of a trade is the return a bank or dealer requires on the capital that regulation makes it
hold against the trade's gross assets, whatever their risk, charged as a running rate on the notional; it makes low-risk, high-volume trades such as
repo-financed bond positions expensive to hold.
\end{definition}

Du, Tepper and Verdelhan found the same force in foreign exchange. Deviations from covered interest parity were large, persistent and not explained by
credit risk or transaction costs. They were strongest for forward contracts that sat on banks' balance sheets at quarter-ends, which points to
regulation as a cause.

\texttt{firm.swapspread} makes the mechanism explicit (\cref{lst:s2:swap-spread-and-asset-swap-trades:model}). For three years the ten-year spread
sits near a 10 basis-point premium. From year 3 a \omterm{def:s2:swap-spread-and-asset-swap-trades:bscost}{balance-sheet cost} of 45 basis points a year phases in over one year, and the spread follows it
down: it averages 11.5 before and $-38.8$ once the cost is in, and it stays negative on every day. At quarter-ends it dips by a further 6 basis points,
12 at year-ends, and a mean-reverting noise of 8 basis points runs throughout (\cref{fig:s2:swap-spread-and-asset-swap-trades:synthetic}).

\begin{omfigure}
\begin{tikzpicture}
\begin{axis}[width=11cm,height=5.4cm,xlabel={\small year},ylabel={\small spread (bp)},tick label style={font=\footnotesize},
  xmin=0,xmax=10,ymin=-75,ymax=30,grid=major,grid style={gray!25},table/col sep=comma,
  legend style={font=\scriptsize,at={(0.97,0.97)},anchor=north east}]
\addplot[omThm,thin] table[x=year,y=spread]{figdata/strategies-2/12-swap-spread-and-asset-swap-trades/synthetic.csv};
\addplot[black,thick,dashed] table[x=year,y=cost]{figdata/strategies-2/12-swap-spread-and-asset-swap-trades/synthetic.csv};
\legend{ten-year swap spread,minus the balance-sheet cost}
\end{axis}
\end{tikzpicture}

\omcaption{The synthetic ten-year swap spread over ten years: a small premium, then a \omterm{def:s2:swap-spread-and-asset-swap-trades:bscost}{balance-sheet cost} phased in from year 3, with quarter-end dips
and noise. Data: \texttt{s2\_swapspread.market}.}\label{fig:s2:swap-spread-and-asset-swap-trades:synthetic}
\end{omfigure}

\section{Trades and their carry}

The long-spread position is one unit of notional: long the ten-year Treasury in repo and paying fixed on a ten-year swap. Its DV01 is 8.5 basis points
of notional per basis point, and the floating rate received exceeds repo by 5 basis points a year. The balance-sheet charge assumes capital of 5\% of
the Treasury notional at a 10\% required return, 50 basis points a year.

\begin{center}
\small
\begin{tabular}{@{}lcccc@{}}
\toprule
bp of notional a year & before, no charge & before, charged & after, no charge & after, charged\\
\midrule
carry (bond yield over swap) & $-11.5$ & $-11.5$ & 38.8 & 38.8\\
marks & $-17.8$ & $-17.8$ & $-11.4$ & $-11.4$\\
floating over repo & 5.0 & 5.0 & 5.0 & 5.0\\
balance-sheet charge & 0.0 & $-50.0$ & 0.0 & $-50.0$\\
total & $-24.3$ & $-74.3$ & 32.4 & $-17.6$\\
Sharpe ratio & $-0.15$ & $-0.44$ & 0.13 & $-0.07$\\
\bottomrule
\end{tabular}
\end{center}

Once the cost is in, the trade has positive carry for anyone who does not pay for balance sheet: 38.8 basis points a year from the negative spread.
For a bank that must hold capital against the Treasury, the charge takes it all and more. The spread settles where the marginal holder is indifferent,
and it stays there because the marginal holder is a bank. The Sharpe ratio is low even without the charge: the spread's noise, marked at a DV01 of
8.5, swamps the carry year to year.

The quarter-end dips are a different trade. Buying the spread at the last close of a quarter and selling five days later earned 103.3 basis points of
notional on average at the five year-ends after the cost was phased in (the worst 90.6), and 43.9 at the other eighteen quarter-ends (the worst
$-4.3$). That is a payment for providing balance sheet on the days regulators measure it. It is planted in the model, and in the real market it is
only as large as the dips are.

\section{Strategy files}

\begin{strategyfile}{Long swap spread}\label{strat:s2:swap-spread-and-asset-swap-trades:long}
\sfield{Who pays you, and why} Banks that cannot hold bonds cheaply; swap receivers (pension funds, insurers) who push swap rates down.
\sfield{Instruments and venues} Treasuries financed in repo; cleared interest rate swaps.
\sfield{Signal} The spread against the holder's own cost of balance sheet.
\sfield{Sizing and execution} DV01-matched; sized on the spread's volatility.
\sfield{Costs} The balance-sheet charge, repo, and swap clearing margin.
\sfield{How it dies} Spreads that go further negative when balance sheet gets dearer.
\sfield{Horizon, capacity, infrastructure} Years; repo lines and swap clearing.
\sfield{Backtest honestly} Charge balance sheet at the holder's real cost.
\sfield{Sources} Jermann (2020); FRED spreads: the thirty-year negative on 93.9\% of days from October 2008 to October 2016; this chapter: 32.4 basis
points a year uncharged, $-17.6$ charged.
\end{strategyfile}

\begin{strategyfile}{Asset-swap carry}\label{strat:s2:swap-spread-and-asset-swap-trades:assetswap}
\sfield{Who pays you, and why} Issuers and investors who leave bonds cheap to swaps.
\sfield{Instruments and venues} Government or agency bonds with a matching swap.
\sfield{Signal} The asset-swap spread against funding cost.
\sfield{Sizing and execution} Bond and swap matched cash flow for cash flow.
\sfield{Costs} Funding and balance sheet.
\sfield{How it dies} The bond's funding cost rises above its asset-swap spread.
\sfield{Horizon, capacity, infrastructure} To maturity or years.
\sfield{Backtest honestly} Funding at the holder's rate, not at the floating index.
\sfield{Sources} No performance figure verified.
\end{strategyfile}

\begin{strategyfile}{Swap-spread curve trade}\label{strat:s2:swap-spread-and-asset-swap-trades:curve}
\sfield{Who pays you, and why} Demand concentrated at one maturity: long-dated receivers at thirty years, bank hedging at shorter ones.
\sfield{Instruments and venues} Spread positions at two maturities.
\sfield{Signal} The spread curve (thirty-year minus ten-year spread) against its history.
\sfield{Sizing and execution} DV01-matched at each maturity.
\sfield{Costs} Four legs.
\sfield{How it dies} A structural change at one end (regulation, pension demand).
\sfield{Horizon, capacity, infrastructure} Months to years.
\sfield{Backtest honestly} The 2008 break: the thirty-year and ten-year spreads moved apart for years.
\sfield{Sources} FRED spreads (this chapter); no performance figure verified.
\end{strategyfile}

\begin{strategyfile}{Year-end balance-sheet trade}\label{strat:s2:swap-spread-and-asset-swap-trades:yearend}
\sfield{Who pays you, and why} Banks that shrink their balance sheets for quarter-end and year-end measurement dates.
\sfield{Instruments and venues} Spread positions, repo, FX swaps across the turn.
\sfield{Signal} The calendar: positions taken at the dip, released after the date.
\sfield{Sizing and execution} Only with balance sheet that is not itself measured on those dates.
\sfield{Costs} Short-dated funding across the turn.
\sfield{How it dies} Rules that average over the quarter rather than measure at its end.
\sfield{Horizon, capacity, infrastructure} Days.
\sfield{Backtest honestly} Measurement dates and rules as they were.
\sfield{Sources} Du, Tepper and Verdelhan (2018) on quarter-end CIP deviations; this chapter's planted dips: 103.3 basis points at year-ends.
\end{strategyfile}

\section{Tutorial: below the Treasury}\label{tut:s2:swap-spread-and-asset-swap-trades:tut}

\begin{tutorial}
\textbf{Goal.} Simulate the swap spread with a \omterm{def:s2:swap-spread-and-asset-swap-trades:bscost}{balance-sheet cost}, run the long-spread trade with and without a charge, and trade the quarter-ends;
read the real spreads. \textbf{End state:} the table and the two figures.
\begin{tutsteps}
\item \textbf{Spread and trade}.
\omcode{code/firm/swapspread/firm_swapspread.py}{44}{74}{The spread with a phased-in balance-sheet cost, and the long-spread trade's parts.}\label{lst:s2:swap-spread-and-asset-swap-trades:model}
\item \textbf{Quarter-ends}.
\omcode{code/strategies-2/12-swap-spread-and-asset-swap-trades/python/s2_swapspread.py}{58}{68}{Buying the quarter-end dip and selling five days later.}\label{lst:s2:swap-spread-and-asset-swap-trades:qe}
\item \textbf{Run} \texttt{s2\_fetch\_swaps.py} once, then \texttt{regimes()}, \texttt{quarter\_ends()} and \texttt{fig\_swapspread.py}.
\end{tutsteps}
\textbf{What to change next.} Make the \omterm{def:s2:swap-spread-and-asset-swap-trades:bscost}{balance-sheet cost} vary with a leverage-ratio rule; add a thirty-year spread with pension demand; run the trade
for an investor whose own \omterm{def:s2:swap-spread-and-asset-swap-trades:bscost}{balance-sheet cost} is 20 basis points.
\end{tutorial}

\section{Build: swap spreads}\label{bld:s2:swap-spread-and-asset-swap-trades:swapspread}

\begin{build}
\bfield{Purpose} A swap spread driven by a \omterm{def:s2:swap-spread-and-asset-swap-trades:bscost}{balance-sheet cost}, the long-spread trade's carry and marks, and a balance-sheet charge.
\bfield{Interface} \texttt{SwapConfig(\dots)}, \texttt{simulate\_spread(cfg)}, \texttt{long\_spread(s, cfg, start, days, charge)}.
\bfield{Rules} Carry is minus the spread; marks are DV01 times the spread's change; the charge is capital times hurdle on the notional.
\bfield{Acceptance tests} \texttt{code/firm/swapspread/tests/}: a flat spread earns its carry exactly; the charge is 50 basis points a year; the
quarter-end dips land on the right days.
\bfield{Stretch} Several maturities; asset swaps on individual bonds; a funding-rate model.
\end{build}

\begin{omsources}
\item U.~J.~Jermann, ``Negative swap spreads and limited arbitrage'', \emph{Review of Financial Studies} 33(1), 2020.
\item W.~Du, A.~Tepper and A.~Verdelhan, ``Deviations from covered interest rate parity'', \emph{Journal of Finance} 73(3), 2018.
\item ICE swap rates and H.15 Treasury yields, via FRED.
\end{omsources}

\section{Exercises}

\begin{exercise}[$\star$]\label{exo:s2:swap-spread-and-asset-swap-trades:1}
A bank holds capital of 5\% against a Treasury position and requires 10\% on capital. What does holding the position cost, in basis points a year?
\end{exercise}

\begin{exercise}[$\star$]\label{exo:s2:swap-spread-and-asset-swap-trades:2}
A position has a DV01 of 8.5 basis points of notional per basis point. The spread widens by 12 basis points. What does the long-spread position gain?
\end{exercise}

\begin{exercise}[$\star$]\label{exo:s2:swap-spread-and-asset-swap-trades:3}
The spread is $-40$ basis points, floating exceeds repo by 5 and the charge is 50. What is the charged carry, a year?
\end{exercise}

\begin{exercise}[$\star\star$]\label{exo:s2:swap-spread-and-asset-swap-trades:4}
Why did the two-year spread stay positive while the thirty-year went negative?
\end{exercise}

\begin{exercise}[$\star\star$]\label{exo:s2:swap-spread-and-asset-swap-trades:5}
Why does the quarter-end dip support the balance-sheet explanation?
\end{exercise}

\begin{exercise}[$\star\star$]\label{exo:s2:swap-spread-and-asset-swap-trades:6}
Who can run the long-spread trade profitably, and why?
\end{exercise}

\begin{exercise}[$\star\star\star$]\label{exo:s2:swap-spread-and-asset-swap-trades:7}
\emph{Coding.} Run the long-spread trade from the start of the sample with no charge. Why does it lose, and what does that say about entering a
carry trade just before its driver changes?
\end{exercise}

\begin{exercise}[$\star\star\star$]\label{exo:s2:swap-spread-and-asset-swap-trades:8}
\emph{Find the flaw.} ``Negative swap spreads are an arbitrage: buy Treasuries, pay fixed, and collect 40 basis points a year risk-free.''
\end{exercise}

\section{Problem: Below the Treasury}

\begin{problem}[{Weekend problem --- swap spreads and balance sheet}]\label{pb:s2:swap-spread-and-asset-swap-trades:1}
The chapter's synthetic spread, FRED's swap rates and the public record.

\medskip\noindent\textbf{Part I --- The spread.}
\begin{enumerate}
\item Define the swap spread and the asset-swap spread.
\item Why were swap spreads positive before 2008?
\item Give the real spreads' averages before and after October 2008.
\item What did Jermann show?
\end{enumerate}

\medskip\noindent\textbf{Part II --- Balance sheet.}
\begin{enumerate}[resume]
\item Define the \omterm{def:s2:swap-spread-and-asset-swap-trades:bscost}{balance-sheet cost}.
\item What did Du, Tepper and Verdelhan find?
\item How does the synthetic spread respond to the cost?
\item Where are the quarter-end dips, and why?
\end{enumerate}

\medskip\noindent\textbf{Part III --- The trades.}
\begin{enumerate}[resume]
\item Define the long-spread trade and its parts.
\item Give its return before and after, with and without a charge.
\item Why is its Sharpe ratio low?
\item What did the quarter-end trade earn?
\end{enumerate}

\medskip\noindent\textbf{Part IV --- The verdict.}
\begin{enumerate}[resume]
\item State the \emph{named result}: the \omterm{def:s2:swap-spread-and-asset-swap-trades:trade}{swap-spread trade}'s return with and without a balance-sheet charge.
\item Why does the spread settle where it does?
\item What would make the spread rise again?
\item How would you size the trade?
\item How would you backtest it honestly?
\item Which strategy file suits a fund with cheap balance sheet?
\item How does this chapter relate to chapter~11?
\item In one sentence: what does a negative swap spread price?
\end{enumerate}
\end{problem}

\section{Interview questions}

\begin{interviewq}[$\star$ \iqroles{trader}]\label{iq:s2:swap-spread-and-asset-swap-trades:1}
What is a swap spread, and what does it mean when it is negative?
\end{interviewq}

\begin{interviewq}[$\star\star$ \iqroles{researcher}]\label{iq:s2:swap-spread-and-asset-swap-trades:2}
Why might the thirty-year swap rate stay below the Treasury yield for years?
\end{interviewq}

\begin{interviewq}[$\star\star$ \iqroles{trader}]\label{iq:s2:swap-spread-and-asset-swap-trades:3}
Build an asset swap on a ten-year Treasury. What do you earn, and what do you pay?
\end{interviewq}

\begin{interviewq}[$\star\star$ \iqroles{risk}]\label{iq:s2:swap-spread-and-asset-swap-trades:4}
A desk is long swap spreads. What scenarios would you stress?
\end{interviewq}

\begin{interviewq}[$\star\star$ \iqroles{developer}]\label{iq:s2:swap-spread-and-asset-swap-trades:5}
How would you charge balance sheet to desks in a firm's P\&L system?
\end{interviewq}

\begin{interviewq}[$\star\star\star$ \iqroles{researcher}]\label{iq:s2:swap-spread-and-asset-swap-trades:6}
Show that the long-spread position's carry is the bond yield minus the swap rate plus the floating rate minus the repo rate, and explain why the
floating-repo gap matters when the swap's floating index changes.
\end{interviewq}
